\begin{abstract}
Models that predict spatial gene expression from H\&E images are usually tested on new patients from
an organ seen during training. This does not tell us whether they work on a new organ. We present
\coast{}, a controlled benchmark for this harder setting. \coast{} separates prediction for new
patients from known organs from transfer to a completely unseen organ. It uses patient-disjoint
splits, selects genes from training data only, and gives every method the same frozen image features,
gene panels, and metrics. Across HEST-1k and STImage-1K4M, all model families perform substantially
worse when the test organ is unseen. We also find that Pearson correlation can become uninformative:
some training-selected genes are almost silent in the test organ, leaving too little target variation
to measure correlation reliably. Finally, we add the same inference-available morphology conditioner
to six representative regression and spatial architectures, including a recent Gene-DML-style
multi-level predictor. It improves all six in both evaluation
settings, with the largest gains for models that do not already combine local and slide-level context.
The gain persists on the recent model and is largest under organ shift. Additional architecture
screening, including neutral and negative cases, is reported in the appendix.
Thus, FiLM is useful across several prediction families, but its benefit depends on the information
already present in a model. \coast{} turns cross-organ generalization from
an implicit claim into a reproducible test. We will release the splits, panels, predictions, and code.
\end{abstract}
